\chapter{Methodology}
\label{chap:methodology}

This chapter describes the methodological decisions, theoretical frameworks, and evaluation strategies underlying the gripper configuration
optimization study. It covers the problem formulation, the choice of optimization library and algorithms, the evaluation architecture, the workpiece
selection rationale, and the comparative experimental design. Implementation details and experimental procedures are treated separately in the
approach chapter.

\section{Research Design}
\label{sec:research_design}

\subsection{Research Paradigm}
\label{subsec:research_paradigm}

The thesis adopts a computational experimental methodology. Both gripper configuration problems are approached through mathematical modeling,
computational-based evaluation, and controlled comparative optimization experiments. This design is appropriate because the objectives governing
gripper quality cannot be expressed in closed analytical form; they depend on geometric interactions between the gripper and the workpiece that must
be assessed computationally. The core methodological contribution is not the development of new optimization algorithms, but the rigorous comparative
evaluation of existing algorithms on well-defined, practically motivated black-box problems.

\subsection{Study Structure}
\label{subsec:study_structure}

The two gripper types are treated as independent optimization studies. The suction-cup gripper optimization seeks a cup layout that simultaneously
maximizes grasp stability and grasp set diversity. The two-finger gripper optimization seeks a jaw contour that achieves high contact quality across a
set of representative workpieces; since the evaluation pipeline computes a separate proximity-based score for each workpiece, each workpiece
contributes one objective, making this a multi-objective problem whose dimensionality scales with the number of workpieces selected. In both cases,
the problem structures and objective spaces differ, and the algorithm comparison is therefore conducted separately for each gripper type with results
not aggregated across them.

\subsection{Experimental Setup}
\label{subsec:experimental_setup}

For each problem, a fixed set of candidate algorithms is applied to a set of representative workpieces. Each algorithm is executed once per workpiece
under identical conditions, including a shared evaluation budget and consistent initialization settings\cite{DaSilva2024Multi-objective}. Algorithm
performance is then analyzed and interpreted across all workpieces and both gripper types, with particular focus on three aspects: convergence
behaviour over the evaluation budget, Pareto front quality as measured by established performance indicators, and the sensitivity of algorithm
rankings to workpiece geometry and objective space dimensionality\cite{Audet2018Performance}. This design allows conclusions to be drawn about which
algorithms are best suited to each problem class and how performance varies with problem instance characteristics\cite{Nikolikj2025Benchmarking}.

\subsection{Scope and Limitations}
\label{subsec:scope_limitations}

Since no repeated independent runs are performed, the study does not support statistical inference about algorithm performance distributions.
Conclusions are therefore descriptive and comparative in nature, grounded in the observed results across the workpiece set rather than in
probabilistic significance testing. Furthermore, the evaluation relies entirely on computational calculation; physical validation of the optimized
gripper configurations is outside the scope of this work.

\section{Problem Formulation}
\label{sec:methodology_problem_formulation}

\subsection{Multi-Objective Optimization Framework}
\label{subsec:methodology_moo}

Both gripper configuration problems are formulated as multi-objective optimization problems. Let $\mathbf{q} \in \mathcal{Q}$
\nomenclature{$\mathbf{q}$}{Decision vector representing a candidate gripper configuration} \nomenclature{$\mathcal{Q}$}{Admissible search space of
decision vectors} denote the decision vector and let $d$
\nomenclature{$d$}{Number of decision variables}
denote the number of decision variables. The general problem is formulated as
\begin{equation}
    \min_{\mathbf{q} \in \mathcal{Q}}\;
    \mathbf{f}(\mathbf{q})
    = \bigl(f_1(\mathbf{q}),\, f_2(\mathbf{q}),\, \dots,\, f_m(\mathbf{q})\bigr),
\end{equation}
where $m$
\nomenclature{$m$}{Number of objective functions}
is the number of objectives, subject to inequality constraints
\begin{equation}
    g_j(\mathbf{q}) \leq 0, \quad j = 1, \dots, n_c,
    \nomenclature{$n_c$}{Number of explicit inequality constraints}
    \nomenclature{$g_j(\mathbf{q})$}{$j$-th inequality constraint function}
\end{equation}
and variable bounds
\begin{equation}
    q_i^{\mathrm{L}} \leq q_i \leq q_i^{\mathrm{U}}, \quad i = 1, \dots, d.
    \nomenclature{$q_i^{\mathrm{L}},\,q_i^{\mathrm{U}}$}{Lower and upper bounds
    of the $i$-th decision variable}
\end{equation}
Objectives that are naturally maximized are negated to maintain a consistent minimization framework throughout. A multi-objective formulation is
preferred over a weighted scalar aggregate because it produces a Pareto-optimal solution set that makes trade-offs explicit and preserves design
alternatives without requiring the prior specification of preference weights.

The number of objectives $m$ is directly determined by the number of workpieces used in the evaluation. This design choice has two methodological
motivations. First, it forces the optimizer to find gripper configurations that perform well across geometrically diverse workpieces simultaneously,
rather than overfitting to a single object. Second, it allows the study to examine algorithm behaviour as the dimensionality of the objective space
increases, since adding workpieces directly increases $m$.

\subsection{Parameterization of Gripper Geometry}
\label{subsec:methodology_decision_variables}

A fundamental methodological decision is how to represent gripper geometry as a finite-dimensional decision vector suitable for black-box optimization
\cite{Ko2020A} \cite{Todescato}. For both gripper types, the representation must satisfy three competing demands: it must be expressive enough to
capture meaningful geometric variation, compact enough to keep the search space tractable, and physically interpretable so that every point in the
search space corresponds to a realizable gripper configuration \cite{Wang2024Fin-Bayes:}.

For the suction-cup gripper, the most general representation would specify the full Cartesian coordinates of each cup independently, yielding a
six-dimensional search space for three cups in the plane. However, this representation introduces redundancy through rigid-body symmetries and admits
degenerate configurations such as collinear or overlapping cups. An alternative is to exploit the geometric structure of the problem: the primary
geometric quantity governing stability is the area of the triangle formed by the three cups, which depends on the radial spread and the relative
angular spacing of the cups rather than their absolute positions. This motivates a reduced polar parameterization that encodes the common radial
distance and the three angular positions, resulting in a four-dimensional search space that eliminates translational and overall rotational redundancy
while retaining the degrees of freedom that directly affect the stability and collision behavior of the layout.

For the two-finger gripper, the finger jaw contour must be represented in a way that allows smooth, physically realizable shapes to be generated from
a small number of parameters \cite{Wolniakowski2017Task}. A spline parameterization is adopted because it provides a compact, smooth, and easily
constrained contour representation, with fixed boundary control points enforcing the geometric constraints at the finger tip and base while the
remaining interior control points form the free decision variables. The number of free control points can therefore be chosen directly to balance
geometric expressiveness against search space dimensionality.

\subsection{Selection of Objective Functions and Evaluation Metrics}
\label{subsec:methodology_objectives}

The choice of objective functions determines what constitutes a good gripper configuration and therefore directly shapes the optimization problem
\cite{Todescato,Wolniakowski2017Task}. For each gripper type, the objectives must reflect physically meaningful performance criteria while remaining
computable from geometry alone, since no physical simulation or hardware trial is available during optimization.

For the suction-cup gripper, two competing criteria govern layout quality. Grasp stability under external disturbances is classically characterized by
force-closure metrics or grasp wrench space volume, both of which require contact force modeling and friction assumptions\cite{Weisz2012Pose}. Since
the optimization operates on geometry alone and no contact force model is available during the search, the area of the support triangle formed by the
three cup positions is used as the stability objective. A larger triangle area directly increases the moment arm resisting tipping disturbances,
making it a geometrically motivated and parameter-free indicator of layout quality\cite{madan2022realrobotchallenge2021}. The second criterion is the
availability of valid grasps on the target workpiece\cite{Bo2023A}. A simple count of collision-free grasp poses does not capture the practical
utility of the grasp set, since many similar poses offer little additional coverage. Diversity-based metrics, which measure the spread of valid grasps
in both position and orientation, are better indicators of operational flexibility and have been used in related work on grasp set
evaluation\cite{Pinskier2024Diversity-Based}. A weighted combination of spatial and orientational diversity is therefore adopted, with both components
normalized to enable consistent comparison across workpieces of different scales.

For the two-finger gripper, the design objective is to find a jaw contour that achieves close contact with the workpiece across as many placements as
possible without collision\cite{Ko2020A}. For the two-finger gripper, the design objective is to find a jaw contour that achieves close contact with
the workpiece across as many placements as possible without collision. The evaluation is performed using a voxel-based convolution scoring function
developed at Fraunhofer IPA, which evaluates finger-to-workpiece proximity over all candidate placements without requiring explicit grasp candidate
enumeration, and excludes colliding placements by construction. This function produces a single scalar score per workpiece that directly serves as the
optimization objective.

\section{Optimization Library Selection}
\label{sec:methodology_library}

All optimization algorithms are implemented and executed using the \texttt{pymoo} framework. This choice is motivated by several methodological
considerations. \texttt{pymoo} is one of the most widely used Python frameworks for multi-objective evolutionary optimization and provides a unified
implementation of a large number of algorithms, which reduces implementation bias and ensures that observed performance differences reflect genuine
algorithmic properties rather than differences in code quality or implementation choices. It includes standardized performance indicators such as
hypervolume, IGD, and GD, as well as built-in constraint handling and visualization tools, which improve the reproducibility and comparability of
results. Furthermore, \texttt{pymoo} is actively maintained, open-source, and thoroughly documented, making it easier for other researchers to
replicate and extend the experiments conducted in this thesis \cite{pymoo}.

\section{Algorithm Selection}
\label{sec:methodology_algorithm_selection}

The selection of optimization algorithms from the set available in \texttt{pymoo} was guided by two criteria: suitability for the problem class and
coverage of fundamentally different optimization paradigms.

Regarding suitability, the gripper configuration problems are continuous, constrained, black-box multi-objective optimization problems. This rules out
a substantial portion of the algorithms available in \texttt{pymoo}. Many algorithms in the library are designed exclusively for single-objective
optimization, including genetic algorithms, differential evolution, particle swarm optimization, CMA-ES, evolution strategies, and direct search
methods, and are therefore not applicable. Several multi-objective algorithms are designed for specialized problem settings outside the scope of this
thesis. These include algorithms for dynamic optimization, where objectives or constraints change over time, and preference-based algorithms, which
require user-defined preference information to guide the search toward specific regions of the Pareto front. Since both gripper problems are static
and the goal is to approximate the full Pareto front without incorporating decision-maker preferences, these specialized methods were
excluded\cite{pymoo}.

Among the remaining general-purpose multi-objective algorithms, the selected methods were chosen because they represent fundamentally different
optimization strategies rather than incremental variations of the same concept:

\begin{itemize}
    \item \textbf{NSGA-II} represents Pareto-dominance-based selection and serves as the standard baseline for multi-objective evolutionary
    optimization\cite{nsga2}.

    \item \textbf{NSGA-III} extends the dominance-based approach through reference directions, making it better suited to problems with a larger
    number of objectives, as arises when more workpieces are included\cite{nsga3}.

    \item \textbf{MOEA/D} represents decomposition-based optimization, solving a set of scalar subproblems rather than relying solely on Pareto
    dominance\cite{moead}.

    \item \textbf{SMS-EMOA} represents indicator-based optimization, directly maximizing the hypervolume indicator and thereby emphasizing Pareto
    front quality\cite{smsemoa}.

    \item \textbf{CMOPSO} provides a swarm-intelligence-based approach, employing a competitive particle swarm mechanism to balance exploration and
    convergence\cite{cmopso}.
\end{itemize}

Together, these five algorithms cover the major paradigms of general-purpose multi-objective optimization, enabling a comparison of fundamentally
different search strategies on the same gripper configuration problems\cite{Doerr2022A, Trivedi2017A, Shang2020A}.

In addition to the evolutionary algorithms, multi-objective Bayesian optimization is included as a comparative strategy. Bayesian optimization is
motivated by its sample efficiency in expensive black-box settings: a probabilistic surrogate model is maintained over the objective functions and an
acquisition function is used to guide evaluations toward promising regions of the search space. Its inclusion allows the study to assess whether
surrogate-based methods offer advantages over population-based methods under the evaluation budgets considered.

\section{Workpiece Selection}
\label{sec:methodology_workpiece_selection}

The workpieces were selected to provide geometric diversity rather than to optimize performance for a particular object. Since the objective of this
thesis is to compare optimization algorithms, it is essential that the benchmark problems exhibit different geometric characteristics and consequently
different optimization landscapes\cite{Cenikj2022SELECTOR:,Yazdani2020Benchmarking}.

The \textbf{IPAIntermediatePlate} represents a simple industrial workpiece with large planar surfaces and regular geometry, reflecting many components
commonly encountered in manufacturing and material handling. Such geometry typically provides a relatively large feasible grasping region and a
smoother objective landscape. In contrast, the \textbf{Stanford Bunny} exhibits a highly irregular three-dimensional shape with varying surface
curvature and complex local features, resulting in a more challenging optimization problem characterized by narrower feasible regions, stronger
collision constraints, and a potentially more multimodal objective landscape\cite{Gupta2020Continuous}.

By evaluating all algorithms on both workpieces simultaneously, the study reduces the risk of drawing conclusions that are specific to a single
geometry and enables a more robust assessment of algorithm performance across problems of varying complexity. The use of multiple workpieces as
separate objectives also means that the comparison implicitly assesses how well each algorithm handles increasing objective space
dimensionality\cite{Cenikj2022SELECTOR:}.

\section{Evaluation Architecture}
\label{sec:methodology_evaluation_architecture}

The evaluation of candidate gripper configurations is decoupled from the optimization process through a client-server architecture. This separation
constitutes a deliberate methodological design choice that enforces a strict black-box interface between the optimization strategy and the objective
function computation. The optimization algorithms operate solely on input-output pairs $(\mathbf{q},\, \mathbf{f}(\mathbf{q}))$ without access to
internal evaluation logic\cite{Bischl2017mlrMBO:}. This architecture also ensures that any optimization algorithm conforming to the interface can be
applied to either evaluation pipeline without modification of either component, supporting reproducibility and extensibility.

\section{Comparative Experimental Design}
\label{sec:methodology_comparative_experiments}

A controlled comparison protocol is adopted to ensure that observed performance differences are attributable to the optimization methods rather than
to differences in experimental conditions. All algorithms are evaluated on the same workpieces, gripper parameterizations, objective functions,
feasibility conditions, and evaluation budget. Each algorithm is executed under multiple independent random seeds to assess not only mean performance
but also variability and robustness across runs\cite{Bouthillier2021Accounting}. Algorithm hyperparameters are initialized from the default values
provided by the \texttt{pymoo} framework \cite{pymoo}. Where the optimization pipeline is configured to do so, a structured hyperparameter tuning
procedure is executed prior to the main optimization run, replacing the defaults with problem-specific settings and reducing the risk of performance
bias from incidental parameter choices. \cite{Beiranvand2017Best}.

\section{Performance Indicators}
\label{sec:methodology_performance_indicators}

Algorithm performance is assessed using standard multi-objective performance indicators. The primary indicator is the \emph{hypervolume indicator},
which measures the volume of objective space dominated by the solution set with respect to a reference point, thereby jointly capturing convergence
quality and Pareto front diversity\cite{Audet2018Performance}. Convergence profiles over function evaluations are also recorded to assess the
efficiency of each algorithm within the given evaluation budget. Robustness is characterized by the distribution of indicator values across repeated
independent runs\cite{Trautmann2009Statistical}.







